% !PRINTED 473
\weberpage{473}
Da, wenn wir uns die Bestimmung von $x$ vorbehalten, die Gleichung (3) jede beliebige gerade Linie darstellen kann, so können wir $x$ so annehmen, dass unter den vier Punkten $y, y', y'', y'''$ keine zwei verbundenen Punkte vorkommen. Setzen wir dann $u_y$ für die unbestimmte Grösse $t$ in (6), so erhalten wir eine rationale Gleichung

\weq{7}{\Psi (u_y, u_x, g, h) = 0,}

und diese Gleichung ist nicht mehr befriedigt, wenn wir für $(x)$ und $(y)$ eine Substitution aus $G_{168}$ machen, durch die $u_x$ geändert und folglich $u_y$ in einen von $u_y, u_{y'}, u_{y''}, u_{y'''}$ verschiedenen Werth übergeführt wird.

Die Gleichung (7) hat also, als Gleichung für $u_x$ betrachtet, nur eine Wurzel mit der allgemeinen Resolvente $R_x$ gemein, und wenn man den grössten gemeinschaftlichen Theiler von $R$ und $\Psi$ aufsucht, so erhält man $u_x$ rational durch $u_y, g, h$ ausgedrückt. Damit ist die allgemeine Resolvente auf die specielle zurückgeführt.

\begin{center}\rule{0.3\linewidth}{0.4pt}\end{center}

Die Gruppe $G_{168}$ enthält noch einen Theiler 21\textsuperscript{sten} Grades $G_{21}$, der durch die Substitutionen $\chi, \tau$ erzeugt wird, und der zu einer Resolvente 8\textsuperscript{ten} Grades Anlass giebt. Als Wurzel dieser Resolvente kann man einfach das Product $x_1 x_2 x_3$ betrachten. Die Co\"efficienten dieser Resolvente sind Invarianten, die bis zum 24\textsuperscript{sten} Grade ansteigen. Wir wollen hier auf diese Resolvente nicht näher eingehen.

\begin{center}
\S.\,127.

\gesperrt{Permutationsgruppe von sieben Ziffern vom Grade 168.}
\end{center}

Da, wie wir gesehen haben, das Formenproblem der ternären Substitutionsgruppe 168\textsuperscript{sten} Grades eine Resolvente 7\textsuperscript{ten} Grades hat, und die Galois'sche Resolvente dieser Gleichung 7\textsuperscript{ten} Grades folglich vom Grade 168 ist, so ergiebt sich daraus, dass in der allgemeinen Permutationsgruppe von sieben Ziffern, deren Grad $1.2.3.4.5.6.7 = 5040 = 168.30$ ist, ein Theiler vom Grade 168 enthalten sein muss. Die Existenz dieses Theilers hat zuerst
%CONT>>
